\section{Introducción}

Los circuitos cuánticos, introducidos por Deutsch en~\cite{Deutsch1989}, son ampliamente reconocidos como un enfoque práctico y de uso común para ilustrar las operaciones de las puertas cuánticas sobre qubits, describiendo en última instancia la matriz unitaria de una computadora cuántica mediante representaciones gráficas.
En consecuencia, los algoritmos y protocolos cuánticos se representan típicamente en forma de circuitos cuánticos.
Esto genera un desafío al tratar con algoritmos o protocolos cuánticos complejos, ya que sus grafos de circuitos pueden superar rápidamente tamaños manejables, haciendo que su representación y manipulación resulten poco prácticas.

En la computación clásica, los circuitos pueden no solo representarse gráficamente, sino también expresarse como expresiones booleanas usando puertas booleanas que implementan lógica basada en tablas de verdad.
Estas expresiones booleanas son muy adecuadas para manipulaciones algebraicas.
Por el contrario, la derivación de una expresión algebraica para un circuito cuántico ha resultado difícil porque la computación cuántica carece de un lenguaje análogo a las expresiones booleanas de la computación clásica.

En~\cite{reina2024ix} Jorge Garcia-Diaz et al. presentan una semántica formal para un novedoso lenguaje algebraico diseñado para
simplificar la expresión de circuitos cuánticos de manera concisa, de modo que las leyes algebraicas fundamentales para circuitos cuánticos puedan demostrarse fácilmente, y los algoritmos cuánticos complejos puedan simplificarse cuando se expresan en este lenguaje.

En este trabajo presentamos una extensión de la notación que nos permite manejar más fácilmente un número arbitrariamente grande de índices; para ello introducimos la notación y las operaciones básicas cuando la indicación de los índices está basada en conjuntos.
